% !TEX root = ../main.tex
% !TEX program = xelatex
\section{模型评价与推广}

\subsection{模型优点与创新}

\begin{enumerate}
    \item \textbf{结构统一，便于逐步扩展。}
    四问共用电量平衡与储能约束，在此基础上逐步加入预测误差、日内调整和波动电价，形成从确定性安排到动态调控的完整方案。预测部分可按信息条件调整，调度部分保持一致，便于解释各问联系和复核运行结果。

    \item \textbf{围绕实际费用处理预测偏差。}
    针对紧急购电价格高的特点，模型根据近期误差设置购电准备，并利用日内实际负载修正后续判断，在提前订电与临时补购之间取得平衡。方案优劣由最终费用和补购风险共同衡量，使预测直接服务于购电决策。

    \item \textbf{区分不同信息带来的收益。}
    在相同模型下，分别检验日内反馈与改单、后发光伏预报的作用，得到前者节省50.63万元、后者额外节省6.95万元。该对照明确了费用改善的来源，为是否引入日内预报提供了直接依据。
\end{enumerate}

\subsection{模型局限与改进方向}

现有参数比较基于2025年数据，跨年适用性仍需更多样本检验。后续可通过跨年滚动验证完善参数校准，并补充供需与电价共同异常的情景分析；实际应用时，再纳入电池老化和网络损耗等成本，使调控方案更贴近长期运行条件。

\subsection{模型推广性评估}

该框架可为校园、社区和工业园区微网提供调控参考。推广时需结合当地供需规律、设备条件和交易规则重新校准参数，并可扩展至风光互补、多台储能及可调负载。通过协调本地发电、储能与外网购电，模型有望兼顾用能经济性与可再生能源利用，为绿色低碳发展提供支持。
